\section{Discussion, limitations, and what would be needed for Clay}\label{sec:discussion}

\paragraph{What this paper achieves.}
This repository and paper deliver a \emph{reproducible decision scaffold} for multiscale reasoning
in the style of the Six Birds framework.
We mechanize (in Lean) the core system-agnostic inequality logic behind a No-Zeno engine
(Section~\ref{sec:no-zeno-engine}), and we mechanize two additional abstract pillars:
(i) route-stability implies polynomial capacity growth (Section~\ref{sec:route-capacity}),
and (ii) a discrete no-go theorem showing that ICAP-style certificates imply pointwise
positive-gain bounds when the feasible input set contains localized tests
(Section~\ref{sec:no-go-icap}).
Separately, we provide a manifest-driven, deterministic Python pipeline that produces
all cited figures and their accompanying metadata (Section~\ref{sec:experiments}).

\paragraph{What this paper does \emph{not} claim.}
We do not prove the Clay 3D Navier--Stokes global regularity theorem.
All PDE-specific hinge claims that would be required for a Clay claim are stated explicitly
as OPEN assumptions elsewhere in the repository. Our emphasis is on (i) making the lemma chain
explicit and auditable, and (ii) showing precisely where and why the ``obvious'' ICAP/capacity
route collapses into a classical peak/strain channel unless one adds a genuine feasibility
restriction that excludes wavepacket-like localization.

\paragraph{Why the ICAP/capacity route does not automatically beat classical criteria.}
A uniform ICAP-style capacity certificate is an \emph{extremal} inequality: it bounds the maximal
positive work that the bridge can extract from \emph{all admissible inputs}.
The discrete no-go theorem (Section~\ref{sec:no-go-icap}) shows that if the feasible input set
is rich enough to contain basis-like localized tests, then any such ICAP bound immediately implies
a pointwise bound on the underlying positive-gain functional. In the PDE interpretation, the
corresponding gain bound aligns with a peak/strain channel (e.g.\ a supremum-type functional
controlling vortex stretching). This explains, structurally, why ``uniform ICAP solves Clay''
is not a free consequence of energy-based accounting.

\paragraph{What we do instead: an ``SBT-legal'' variant.}
Given the above obstruction, the only viable bypass within this scaffold is to restrict feasibility
so that localized wavepacket tests are \emph{not} admissible. We formalize this as an anti-localization
hinge and show that it can follow from finite-channel delocalized feasibility in a discrete model.
We then propose an operational certificate (computed by scripts in this repository) that checks
anti-localization-style metrics for initial data (Section~\ref{sec:spt-legal}).
This does not claim that legality is preserved under Navier--Stokes evolution; rather, it proposes
a concrete restricted theory where the No-Zeno engine becomes meaningful without smuggling in a
classical supremum regularity hypothesis.

\subsection{OPEN gaps: what would be needed for a Clay claim}
To turn this scaffold into a Clay regularity proof, one would need to establish PDE-facing statements
that go beyond what is currently proved here. Concretely:
\begin{itemize}
\item \textbf{A channelization theorem for the canonical NS packaged bridge:}
prove (or refute) that the feasible interface at shell $j$ is effectively low-channel and delocalized
in a sense strong enough to exclude wavepacket-like tests.
\item \textbf{Anti-localization as a theorem (not a certificate):}
prove that the actual NS-relevant feasible inputs satisfy an anti-localization inequality with
$\eta_j\to 0$ (or prove that intermittency produces counterexamples).
\item \textbf{A PDE bridge from ``No-Zeno'' to smooth continuation:}
show that the absence of a Zeno cascade (as defined by a canonical frontier) yields a continuation
criterion in an $H^s$ class, closing the loop from abstract increments to PDE regularity.
\end{itemize}

\subsection{Future work}
\paragraph{Theory (paper-facing).}
\begin{itemize}
\item Prove or refute a concrete Navier--Stokes statement implying wavepacket-exclusion for feasible inputs
(e.g.\ channelization/delocalization for a canonical packaged bridge object).
\item Clarify whether any such wavepacket-exclusion principle can be derived from NS structure without implicitly
assuming a classical peak-channel bound (BKM/Serrin-type) \citep{bkm_1984,serrin1962interior}.
\item Mechanize additional abstract ``interface calculus'' steps in Lean as they are used in the scaffold,
including alternative discrete feasibility models and their implications.
\end{itemize}

\paragraph{Numerics (repo-facing).}
\begin{itemize}
\item Push the existing diagnostic suite to higher $N$, longer horizons, and broader forcing/viscosity regimes,
and quantify robustness of channel and anti-localization metrics.
\item Stress-test intentionally illegal constructions (multiple localized patterns, multi-shell localization,
intermittency-like bursts) to map how the certificate fails.
\item Add sensitivity analyses: dependence of $\kappa_j$ and effective ranks on estimator choices (window sizes,
pooling strategy, basis choices).
\end{itemize}
